\documentclass[11pt]{article}
\usepackage{amsmath,amssymb,amsthm,booktabs}
\usepackage[margin=2.5cm]{geometry}
\newtheorem{theorem}{Theorem}\newtheorem{lemma}[theorem]{Lemma}\newtheorem{proposition}[theorem]{Proposition}
\newtheorem{corollary}[theorem]{Corollary}\newtheorem{conjecture}[theorem]{Conjecture}
\theoremstyle{definition}\newtheorem{definition}[theorem]{Definition}\newtheorem{remark}[theorem]{Remark}
\title{Q597: the second journey across one star\\
\normalsize PROVE 2026-11-18 --- one-star SCEL \emph{not proved}; reduced to a journey-level block-swap hypothesis (BS);\\
``BS $\Rightarrow$ local formula, one side only'' \emph{proved}; BS \emph{computed} (521/521, control 0/880)}
\author{Clio}\date{18 November 2026}
\begin{document}\maketitle

\section*{Provenance}
\begin{itemize}
\item Code: \texttt{proofs/code-1118-Q597/} (new; \texttt{ls} before \texttt{mkdir} showed no \texttt{code-1118*}). Scripts: \texttt{widen.py}, \texttt{widen6.py} (step 0),
\texttt{align.py} (records every Purbhoo step of all four journeys per pair, \texttt{align.pkl}), \texttt{analyse\_align.py}, \texttt{middle.py}, \texttt{words.py},
\texttt{transl.py}, \texttt{stepdelta.py}, \texttt{reduce.py}, \texttt{messy.py}, \texttt{local.py}, \texttt{lesion.py}, \texttt{waketr.py}, \texttt{oneside.py}; helpers in \texttt{q597lib.py}.
The engine is unchanged (\texttt{code-1112c2-Q585/migrate.py}); steps are recorded by wrapping \texttt{Board.step}. All arithmetic is exact (integers and \texttt{Fraction}s in $\mathbb Z^4$).
\item Read first-hand: \texttt{migrate.py} (all of \texttt{Board.step}, \texttt{journey}), Purbhoo 0705.1184 \S3.1 (migration) and \S4 (wakes, Wake Crossing Lemma) from the seed PDF,
\texttt{2026-11-17c2-MC-pq-exchange.tex}, \texttt{2026-11-17-PRIGID.tex} \S Setting, and today's browse log.
\item \textbf{Not read:} Fulton \emph{Young Tableaux} \S A.3 Lemma 2 and Brown--Moran 2503.11840 (the browse pointer). The session rules forbid browsing and neither is on disk.
Purbhoo \S4 turned out to be the closer template, because it is about the same objects.
\end{itemize}

\section*{Summary}
\begin{enumerate}
\item \textbf{Step 0.} The 11-17 cap was \emph{not binding}: caps 60 and 300 both visit exactly 5194 states. So the 11-17 ``sample'' is the \emph{entire} reachable set for the census
$|\lambda|,|\mu|\le5$, $\mathfrak{gl}_4$. SCEL with the 11-17 signatures frozen: 1558/1558; on-$I$ law 1558/1558; $B/B$ control 440 non-commuting, rule predicts none.
A genuine widening ($|\lambda|$ or $|\mu|=6$) was launched; its result is in \S\ref{sec:step0}.
\item \textbf{Step 1.} \texttt{Board.step} \emph{is} Purbhoo's step (minimal hexagon weakly right of the rhombus; rotate by $180^\circ$ or $120^\circ$). It reads a window of tiles within
three edge-adjacencies of the rhombus, and writes exactly the four tiles of one hexagon (Lemma~\ref{lem:window}). Each step changes \emph{exactly two} pure cells and keeps
their keys (computed: 31\,809 steps). A journey's pure-cell delta is the sum of its step deltas (4674/4674 journeys).
\item \textbf{Step 2, computed.} Align $y$'s journey in $M$ with $y$'s journey in $Mx$. On \emph{every} $B/C$ pair, the two journeys agree on a prefix and a suffix and differ on a
\emph{middle} of equal length (2--9 steps), after which the rhombus is in the \emph{same} tile (re-synchronisation, 1558/1558 records $\times$ both journeys).
On the 440 non-commuting $B/B$ pairs the middles have unequal multisets and never re-synchronise (0/880).
\item \textbf{Block swap (BS), computed.} On the interacting side's ``clean'' journey ($y$ if the interaction is on $q$, $x$ if on $p$), the middle in $M$ is $a\cdot C$, and in the other
state it is $(C-\alpha)\cdot(a+\gamma)$. Here every step is a rigid $\mathbb Z^4$-translate of the corresponding step, $\alpha$ is the rhombus displacement of $a$, and $\gamma$ that of $C$:
253/253 ($q$) and 268/268 ($p$). The rhombus goes round a parallelogram with sides $\alpha,\gamma$. This is the rhombus 2-cell of 11-17, one level up, at the level of journeys.
\item \textbf{Proved:} BS $\Rightarrow$ an exact formula for the second-move change (Theorem~\ref{thm:red}). It is supported on a bounded patch and lives on \emph{one side only}
(Corollary~\ref{cor:side}). The formula reproduces the observed change 521/521.
\item \textbf{Not proved:} BS itself, i.e.\ the crossing analysis. Nor (E): the middles do not read every cell of $I$, so (E) needs window data. Nor the other journey: it has no translate
structure, and Theorem~\ref{thm:red} only covers one journey per side, so SCEL $\Rightarrow$ (MC) still needs it. Precise gaps in \S\ref{sec:gaps}.
\end{enumerate}

\section{Setting}
As in \texttt{2026-11-17c2-MC-pq-exchange.tex}. A state $M$ is a mosaic (a twisted fine mixed subdivision in $\mathcal F$ of PRIGID.tex, Def.~\texttt{def:fms}). Its tiles are
\texttt{tri} (pure cells), \texttt{sq} (mixed, edges at $90^\circ$) and \texttt{rh} (mixed, edges at $30^\circ$; outside nests only during a journey). Vertices lie in
$\mathbb Z^4=\Lambda_0\oplus\Lambda_1$. A pure cell has key $(0,\tau)$ with value $q(\tau)\in\Lambda_1$, or key $(1,\sigma)$ with value $p(\sigma)\in\Lambda_0$. For a move $m$,
$\Delta_M(m)$ is the net change of values (11-17c2). For a pair $(x\to B,\,y\to C)$: $X=\Delta_M(x)$, $Y=\Delta_M(y)$, $Y^x=\Delta_{Mx}(y)$, $X^y=\Delta_{My}(x)$.
A \emph{journey} is the sequence of steps of one rhombus. A \emph{step} is a triple $\sigma=(\mathrm{typ},g,n)$: the set $g$ of four tiles it removes and the set $n$ of four it adds.
For $v\in\mathbb Z^4$, $\sigma+v$ translates every vertex. On deltas, $(T_v\delta)(i,\kappa+v_i)=\delta(i,\kappa)$, where $v_0,v_1$ are the two blocks of $v$.

\section{Step 0: the census is exhausted, then widened}\label{sec:step0}
\begin{center}\begin{tabular}{lrrrrr}\toprule
run & states & $B/C$ & SCEL $I$ & on-$I$ law & $B/B$ non-comm.\ (pred.\ none)\\\midrule
cap 60 (11-17 setting) & 5194 & 1558 & 1558 & 1558 & 440\\
cap 300 & 5194 & 1558 & 1558 & 1558 & 440\\
\multicolumn{6}{l}{\emph{$|\lambda|$ or $|\mu|=6$ (new census entries, uncapped):} see the addendum at the end.}\\\bottomrule
\end{tabular}\end{center}
$|U|$ histogram (both caps): $q$: 143, 79, 30, 1 for $|U|=1,2,3,4$; $p$: 219, 49 for $|U|=1,2$; empty 1037.
Lesson for the record: the 11-17 brief's worry (cap 60 skips states) was right to check, and the answer is that nothing was skipped. The real limit is the census, not the cap.

\section{Step 1: what one step reads and writes}
\begin{lemma}[window]\label{lem:window}
Let $T$ be a board, $R$ the migrating rhombus, $d$ the direction. Let $N(R,T)$ be the set of tiles reachable from $R$ through at most three edge-adjacencies, the first
three through non-rhombus tiles. Then $\mathrm{step}(T,R)$ is a function of $(R,N(R,T),d)$ alone, and it removes and adds exactly four tiles: $R$ and three tiles of $N(R,T)$.
\end{lemma}
\begin{proof}
From the definition (Purbhoo \S3.1; \texttt{Board.step}, read line by line). The candidates are the connected sets $S\ni R$ grown three times by adjoining an edge-neighbour
that is not a rhombus, keeping only those with tile multiset $\{\mathrm{rh},\mathrm{sq},\mathrm{tri},\mathrm{tri}\}$. So every candidate lies in $N(R,T)$. Each test applied to a candidate
reads only the vertices of its tiles and of $R$: its boundary cycle has six corners, it lies weakly right of $\min x(R)$, and it is symmetric under rotation by $180^\circ$ or $120^\circ$ about its vertex centroid.
The tie-break reads the tile across one edge of $R$, which lies in $N$. The output rotates the four tiles of the chosen $S$ and nothing else.
\end{proof}
\begin{lemma}[translation covariance]\label{lem:cov}
If $N(R+v,T')=N(R,T)+v$ then $\mathrm{step}(T',R+v)=\mathrm{step}(T,R)+v$.
\end{lemma}
\begin{proof}
The projection $\varphi$ is linear, so each test of Lemma~\ref{lem:window} compares differences, and these are unchanged by $+v$. ``Weakly right of $\min x(R)$'' shifts both sides by $x(\varphi v)$.
Rotation about the vertex centroid commutes with translation, and the tie-break is a statement about edges of $R$. Hence the candidate sets, the chosen $S$ and the rotation translate by $v$.
\end{proof}
\begin{lemma}[two-cell steps; additivity]\label{lem:add}
Assume every board met during a journey lies in $\mathcal F$, so that each key occurs exactly once (PRIGID.tex, Def.~\texttt{def:fms}). Then every step preserves the set of pure-cell keys
inside its hexagon, so it has a well-defined delta $\delta_\sigma$. The net delta of a journey $\sigma_1\cdots\sigma_N$ is $\sum_i\delta_{\sigma_i}$, and $\delta_{\sigma+v}=T_v\delta_\sigma$.
\end{lemma}
\begin{proof}
A step changes tiles only inside its hexagon. Keys outside are untouched, and by the exactly-once property the keys inside are the same set before and after. So the per-step delta is defined,
and the deltas add along the journey (telescoping). Translating a tri tile $\tau\times\{w\}$ by $v$ gives $(\tau+v_0)\times\{w+v_1\}$, and the change of value is unchanged.
\end{proof}
\noindent\emph{Computed:} on all 31\,809 recorded steps the key set is preserved, and each step changes exactly two pure cells. Type a changes two triangles of one side: 13\,569 on side 0, 9590 on side 1.
Type b changes one triangle of each side: 8650. Additivity holds on 4674/4674 journeys. The hypothesis ``intermediate boards lie in $\mathcal F$'' is \emph{not} certified. PRIGID control (i)
covers mosaics at rest, not mid-journey boards. It is gap G5.

\begin{proposition}[prefix and suffix]\label{prop:prefix}
Let $T,T'$ be boards with the same rhombus $R$, and let $D=T\triangle T'$. If the windows $N(R_k,T_k)$ of the journey $\sigma_1,\sigma_2,\dots$ in $T$ avoid the current $D$ for
$k\le m$, and the termination test agrees for $k\le m$, then the journey in $T'$ begins with the same steps $\sigma_1,\dots,\sigma_m$, and $D$ is unchanged after them.
\end{proposition}
\begin{proof}
Induction on $k$. If $N(R_k,T_k)\cap D=\emptyset$ then $N(R_k,T'_k)=N(R_k,T_k)$: tiles are polygons of a tiling, and the window is generated from $R_k$ by adjacency inside tiles common to both.
Lemma~\ref{lem:window} gives the same step. That step removes and adds the same tiles in both boards, and all of them lie outside $D$, so $D$ is unchanged.
\end{proof}
The same statement read from the end gives the suffix. Once the rhombi coincide again and the remaining windows avoid the new difference, the journeys finish identically.
\emph{Caveat:} the termination test (\texttt{journey}: is the rhombus packed into a nest $\ne$ source?) reads nest data, and nest $B$ differs between $M$ and $Mx$. In all records
$y$ lands in nest $C$ in both boards.

\section{Step 2: the journeys re-synchronise}
For each pair, write the journey in $M$ as $P\cdot m\cdot S$ and in the other state as $P\cdot m'\cdot S$, with $P$ the longest common prefix and $S$ the longest common suffix (\texttt{analyse\_align.py}).
\begin{center}\small\begin{tabular}{lrll}\toprule
class & pairs & $(|m|,|m'|)$, $y$-journey & $(|m|,|m'|)$, $x$-journey\\\midrule
$B/C$ disjoint wakes & 779 & $(0,0)$ & $(0,0)$\\
$B/C$ overlapping, $I=\emptyset$ & 258 & $(2,2),(3,3)$ & $(2,2),(3,3)$\\
$B/C$, $q$, $|U|=1$ & 143 & $(3,3),(4,4)$ & $(4,4),(5,5)$\\
$B/C$, $p$, $|U|=1$ & 219 & $(4,4)$--$(7,7)$ & $(3,3)$--$(6,6)$\\
$B/C$, $|U|\ge2$ & 159 & $(4,4)$--$(9,9)$ & $(4,4)$--$(8,8)$\\
$B/B$ commuting & 339 & $(0,0)$ & $(0,0)$\\
$B/B$ non-commuting & 440 & unequal, e.g.\ $(4,6)$ & unequal, e.g.\ $(4,3)$\\\bottomrule
\end{tabular}\end{center}
\textbf{Re-synchronisation (computed):} in every $B/C$ record, for both journeys, the rhombus after $m$ equals the rhombus after $m'$: 1558/1558 per journey.
On $B/B$ non-commuting pairs it never does (0/440 per journey).

\paragraph{Block swap.} Let the \emph{clean journey} be $y$ (second in $Mx$) when the interaction is on side $q$, and $x$ (second in $My$) when it is on side $p$.
\begin{center}\small\begin{tabular}{lrl}\toprule
& pairs & structure of $(m,m')$ \\\midrule
$q$, clean ($y$) & 253 & $m=a\cdot C$, $m'=(C-\alpha)\cdot(a+\gamma)$, all steps rigid translates \\
$p$, clean ($x$) & 268 & same \\
$q$, other ($x$) / $p$, other ($y$) & 253 / 268 & 11 / 39 rotations; the rest have steps of shapes absent from $m$ \\
$I=\emptyset$, overlapping & 258 & same as the clean form on 218 ($y$) and 244 ($x$) \\
$B/B$ non-commuting & 440 & never (0/880)\\\bottomrule
\end{tabular}\end{center}
Here $\alpha$ is the rhombus displacement of the first middle step, and it is \emph{one vector per side}: $\alpha=(0,-1;0,0)$ on $q$ (253/253) and $\alpha=(0,0;0,1)$ on $p$ (268/268).
$\gamma$ is the total displacement of $C$. Example ($q$, \texttt{middle.py}): $m=(\mathrm a,\mathrm b,\mathrm b)$, $m'=(\mathrm b,\mathrm b,\mathrm a)$. The rhombus goes
$\rho\to\rho+\alpha\to\rho+\alpha+\gamma$ in $M$ and $\rho\to\rho+\gamma\to\rho+\gamma+\alpha$ in $Mx$, and every tile of the crossing is shifted by exactly $-\alpha$.
It is a commuting square of journeys, and it is the 11-17 rhombus $\{v,v+b,v+s,v+c\}$ one level up.

\begin{definition}[BS]
A pair satisfies \emph{BS} if its clean journey decomposes as $P\cdot a\cdot C\cdot S$ in $M$ and $P\cdot(C-\alpha)\cdot(a+\gamma)\cdot S$ in the other state.
\end{definition}

\section{Step 2, proved part: BS forces a local, one-sided change}
\begin{theorem}[reduction]\label{thm:red}
Assume Lemma~\ref{lem:add}'s hypothesis. If BS holds for the clean journey (say $y$; the other side is symmetric), then
\[Y^x-Y\;=\;\sum_{\sigma\in C}\big(T_{-\alpha}\delta_\sigma-\delta_\sigma\big)\;+\;T_{\gamma}\delta_a-\delta_a .\]
In particular $Y^x=Y$ off $K=\mathrm{keys}(a)\cup(\mathrm{keys}(a)+\gamma)\cup\mathrm{keys}(C)\cup(\mathrm{keys}(C)-\alpha)$, a set of at most $4(|C|+1)$ cells.
\end{theorem}
\begin{proof}
By Lemma~\ref{lem:add}, $Y=\sum_P\delta+\delta_a+\sum_C\delta+\sum_S\delta$ and $Y^x=\sum_P\delta+\sum_{C}\delta_{\sigma-\alpha}+\delta_{a+\gamma}+\sum_S\delta$. The common terms cancel.
Then $\delta_{\sigma+v}=T_v\delta_\sigma$. Every $\delta_\sigma$ is supported on the keys of $\sigma$, so the right side is supported on $K$.
\end{proof}
\begin{corollary}[one side only]\label{cor:side}
In Theorem~\ref{thm:red}, suppose $\alpha$ lies in a single block $\Lambda_i$ and $\delta_a$ is supported on side $i$. Then $Y^x-Y$ is supported on side $i$.
Crossing steps whose delta lives on side $1-i$ contribute zero, and a type-b step contributes only through its side-$i$ cell.
\end{corollary}
\begin{proof}
$T_{-\alpha}$ shifts side-$(1-i)$ keys by the $(1-i)$-block of $-\alpha$, which is $0$, so $T_{-\alpha}\delta_\sigma-\delta_\sigma$ vanishes on side $1-i$. The term $T_\gamma\delta_a-\delta_a$ is supported
on side $i$ because $\delta_a$ is.
\end{proof}
\noindent\textbf{Computed checks of the hypotheses and conclusions} (\texttt{reduce.py}, \texttt{oneside.py}, \texttt{local.py}):
\begin{itemize}
\item The formula of Theorem~\ref{thm:red} equals the observed $Y^x-Y$ (or $X^y-X$ on $p$) in \textbf{521/521} records. Its support is $4|U|+2$ cells: $6,10,14,18$.
\item $\alpha$ lies in one block and $\delta_a$ is on that side: 253/253 ($q$, block $\Lambda_0$, side $q$) and 268/268 ($p$, block $\Lambda_1$, side $p$). This \emph{explains} the 11-17 observation
that the interaction never occurs on both sides, given BS.
\item Crossing steps that contribute: exactly $|U|+1$ in every record, on both sides. The remaining crossing steps are other-side type-a steps with zero contribution: 35 ($q$), 180 ($p$).
On $p$ the crossing can be long for a single star ($C=v_1v_3^kv_1$, $k\le3$), because $y$ runs along $x$'s wake. Corollary~\ref{cor:side} is why the support stays six cells.
\item Prefix and suffix never touch $I$ (521/521), so $Y^x=Y$ there is not a cancellation accident.
\item Only \textbf{8} step shapes up to translation occur in clean middles. The local configurations $(m,m')$ up to translation number 2 ($q$, $|U|=1$), 3, 3, 1 ($q$, $|U|=2,3,4$), 4 ($p$, $|U|=1$)
and 10 ($p$, $|U|=2$).
\end{itemize}
\paragraph{Consequence for (L).} For the clean journey, (L) (no change off $\mathrm{star}(u)$) follows from BS plus the finite check that, for each of the 2 ($q$) and 4 ($p$) one-star
configurations, the formula is supported on one star. The check passes on all records. For the brief's literal (L), about $y$ on \emph{both} sides, the $p$-side needs the other
journey, which is not covered (G3).

\section{Step 3: (E)}
(E) on $I$: $X+Y^x=s$ (11-17c2), computed 521/521. \textbf{Lesion:} with $b$ or $c$ in place of $s$ it fails on 521/521 records each (\texttt{lesion.py}).
It does \emph{not} follow from Theorem~\ref{thm:red} alone. $X$ on $I$ is $x$'s solo delta, and the middles do not read every cell of $I$: in 521/521 records some cell of $I$ is read by only
one of $m,m'$. For instance, $U_0$ has signature $(0,0)$ and is touched only in $m'$. So (E) needs the values of $M$ on star cells that $y$ never moves, i.e.\ the windows of the middle steps
and not just their tile sets. The brief's Step 3 (local lemma over realisable star configurations from \texttt{mixsub.py}) was \emph{not run}. It is gap G2.

\section{A refuted route}
\emph{Hope:} a journey acts on its wake by translating each wake tile, which would make BS a corollary of Lemma~\ref{lem:cov}.
\emph{False} (\texttt{waketr.py}): in 1520 of 3116 solo journeys some removed tile is not a translate of any added tile. A $180^\circ$ step flips a triangle's orientation.
So BS is a statement about the \emph{windows of the crossing}, not about whole wakes.

\section{Gaps}\label{sec:gaps}
\begin{description}
\item[G1 (BS, the core).] Prove that when the clean journey's window first meets $D=M\triangle Mx$, the board around it is the translate by $-\alpha$ of the board that the crossing $C$ meets
in $M$ after $a$, and that after $C-\alpha$ the step $a+\gamma$ re-joins. This is a crossing analysis in the style of Purbhoo's Wake Crossing Lemma (\S4, cases (a)--(f)), but for journeys
towards \emph{different} nests. Purbhoo's lemma is about rhombi of one flock that never cross a midline; here $y$ \emph{does} cross $x$'s wake. Computed 521/521; control 0/880.
\item[G2 ((E)).] It needs window data at the star. Proposed: extend each local configuration by the windows of its steps and verify (E) per configuration. That is finite if G4 holds.
\item[G3 (the other journey).] $x$ in $My$ ($q$) and $y$ in $Mx$ ($p$) re-synchronise (1558/1558), but have no translate structure. Theorem \texttt{thm:cond} of 11-17c2 needs both $Y^x$ and $X^y$.
So even BS + G2 give only half of SCEL. A possible route is Purbhoo's Theorem~2 (invertibility) with a mirror symmetry exchanging $B\leftrightarrow C$, $0\leftrightarrow1$. Untested.
\item[G4 (completeness).] The catalogue (2+4 one-star configurations) is what occurs in the census, not a proved list.
\item[G5.] Intermediate boards in $\mathcal F$ (hypothesis of Lemma~\ref{lem:add}). Computed on 31\,809 steps.
\end{description}

\section{Status}
\begin{itemize}
\item \textbf{Proved:} Lemmas~\ref{lem:window}, \ref{lem:cov}, \ref{lem:add} (under G5); Proposition~\ref{prop:prefix}; Theorem~\ref{thm:red}; Corollary~\ref{cor:side}.
\item \textbf{Computed:} re-synchronisation (all $B/C$; $B/B$ 0/440); BS 521/521 ($B/B$ 0/880); formula 521/521; (E) 521/521 with lesion 0/521; census exhausted at caps 60/300.
\item \textbf{Not proved:} one-star SCEL ((L) and (E)). Hence (MC) remains computed. The retreat (Q596(b)) was not run.
\end{itemize}

\end{document}
